\section{Referencias}

\begin{itemize}
    \item \textbf{OpenMP Application Programming Interface Examples, Version 6.0.1 (noviembre 2025)} --- fuente oficial de los ejemplos utilizados para justificar la elección de las directivas \texttt{parallel for collapse(2)} (secciones 3.1, 3.8 y 5.3.5). Disponible en \texttt{github.com/OpenMP/Examples} (etiqueta \texttt{v6.0.1}).

    \item \textbf{OpenMP Application Programming Interface Specification, Version 6.0 (noviembre 2024)} --- especificación del estándar que define directivas, cláusulas, rutinas de biblioteca y variables de ambiente. Disponible en \texttt{openmp.org}.

    \item \textbf{GNU \texttt{gprof}} --- perfilador por muestreo con instrumentación en tiempo de compilación (\texttt{-pg}). Manual en \texttt{sourceware.org/binutils/docs}.

    \item \textbf{Valgrind: \texttt{cachegrind}} --- simulador de caché del proyecto Valgrind. Documentación en \texttt{valgrind.org/docs/manual/cg-manual.html}.

    \item \textbf{POSIX.1-2017 (IEEE Std 1003.1-2017)} --- \texttt{clock\_gettime}, \texttt{getrusage}, \texttt{sched\_getaffinity}. Especificación de la API POSIX utilizada para la medición del tiempo y de la memoria residente.

    \item \textbf{Intel\textsuperscript{\textregistered} 64 and IA-32 Architectures Software Developer's Manual, Volume 2} --- referencia de los contadores de rendimiento del procesador y del modelo de jerarquía de caché.

    \item \textbf{Brendan Gregg --- \textit{Systems Performance}} y su blog \texttt{brendangregg.com} --- referencias sobre las herramientas de línea de comandos para rendimiento en Linux.

    \item \textbf{Material del curso de Computación de Alto Rendimiento} --- presentaciones y grabaciones de clase (OpenMP, benchmark, profiling) utilizadas como referencia conceptual durante el desarrollo del caso de estudio 2.
\end{itemize}
